\documentclass[12pt, a4paper]{article}

% Language
\usepackage[utf8x]{inputenc}
\usepackage[T1]{fontenc}
\usepackage[spanish]{babel}

\usepackage{amsmath, amssymb, amsthm, amsfonts, mathrsfs}
\usepackage[]{graphicx}
\usepackage[table,xcdraw]{xcolor}
\usepackage{booktabs}
\usepackage{longtable}
\usepackage[font={small}]{caption}
\usepackage{multirow}
\usepackage{float}
\usepackage{tikz}
\usepackage{minted2}
\usepackage{verbatim}
\usepackage{caption}
\usepackage{subcaption}
\usepackage{gensymb}

\usepackage{hyperref}
\hypersetup{
  colorlinks = true,
  linkcolor = black,
  urlcolor = blue,
  citecolor = black}

% Margins
\usepackage[
  a4paper,
  includehead,
  footskip=7mm, headsep=6mm, headheight=7mm,
  top=20mm, bottom=20mm, left=25mm, right=25mm
]{geometry}

% Paragraph
\setlength{\parindent}{0cm} % No indent
\setlength{\parskip}{6pt}
\linespread{1.15}
\hyphenpenalty=0

% Header/Footer
\usepackage{fancyhdr}
\fancyhf{}
\setlength{\headheight}{20pt}
\addtolength{\topmargin}{-2pt}
\fancyhead[L,R]{\bfseries\thepage}
\fancyhead[L]{Diseño e implementación de un lenguaje}
\fancyhead[R]{Especificación}
\fancyfoot[R]{\thepage}
\renewcommand{\headrulewidth}{0.5pt}
\renewcommand{\footrulewidth}{0.5pt}
\fancyheadoffset[R,L]{0\textwidth}
\pagestyle{fancy}

% \usetikzlibrary {graphs,graphdrawing} \usegdlibrary {trees}


\begin{document}

\begin{titlepage}
  \begin{center}
    \vspace*{1cm}
    \textbf{Autómatas, Teoría de Lenguajes y Compiladores}

    \vspace{1cm}
    \large{Diseño e implementación de un lenguaje}\\
    \Large\textbf{Especificación}

    \vspace{1cm}
    \large\textbf{Teo Fumagalli}

    \vfill

    \normalsize
    \textbf{Ingeniería Informática, ITBA}\\
    \textbf{Buenos Aires, Argentina}\\
    \textbf{\today}
    \vspace*{1cm}

  \end{center}
\end{titlepage}


\tableofcontents
\pagebreak

\section{Equipo}

\begin{table}[H]
	\centering
	\begin{tabular}{|l|l|l|l|}
		\hline
		\textbf{Nombre} & \textbf{Apellido} & \textbf{Legajo} & \textbf{E-mail} \\
		\hline
		Teo             & Fumagalli         & 66048           & \href{mailto:tfumagalli@itba.edu.ar}{\texttt{tfumagalli@itba.edu.ar}} \\
		\hline
	\end{tabular}
\end{table}

\section{Repositorio}

La solución será versionada en \href{https://github.com/teofum/tpe-tla}{https://github.com/teofum/tpe-tla}

\section{Dominio}

Desarrollar un lenguaje para generación de sitios web estáticos (HTML + CSS +
JS) que puedan ser hosteados por un servidor web estándar (e.g. Apache, nginx).
El lenguaje debe permitir crear programas cuya salida sea una página HTML
válida, crear componentes reutilizables en la forma de funciones, y programar
lógica sencilla (condicionales, iterar listas), y leer datos de entrada de
ciertos formatos (JSON, CSV) a variables.

Con tal de reducir la complejidad del proyecto, los estilos (CSS) y cualquier
comportamiento frontend (JS) serán importados como \textit{assets} externos y
referenciados en el HTML generado.

Toda la evaluación del programa que genera una página HTML será realizada dentro
del compilador como parte del proceso de compilación. El artefacto de entrada es
uno o más archivos fuente o un directorio con archivos fuente y recursos (CSS,
JS, imágenes, datos en formato JSON o CSV, etc.), y el artefacto de salida es un
directorio con archivos HTML correspondiendo a cada archivo fuente que retorne
una página completa, y todos los recursos necesarios.

La implementación satisfactoria del lenguaje permitiría escribir sitios web
sencillos (estáticos) con un \textit{toolchain} muy simple y sin dependencias
externas (e.g. npm), ningún componente requerido en tiempo de ejecución, y en un
lenguaje más amigable que HTML estándar.

\section{Construcciones}

El lenguaje desarrollado debería ofrecer las siguientes construcciones,
prestaciones y funcionalidades:

\subsection{Variables}

El lenguaje permitirá definir variables con un \textit{scope} localizado al
bloque más cercano (con reglas similares a e.g. C). Cada variable tiene un tipo
y un valor; las variables deberán ser inicializadas explícitamente al definirlas
y no hacerlo será un error.

También se podrán definir constantes, con semántica idéntica a una variable con
la excepción de que no pueden ser modificadas una vez definidas. Se permitirá la
definición de variables con el mismo nombre que otras (\textit{shadowing}), pero
nunca dentro del mismo \textit{scope}.

\paragraph{Ejemplos de sintaxis}
\small
\begin{verbatim}
x: int = 5           // Declara la variable "x" con tipo "int" y valor 5
y := 5               // Se infiere el tipo "int" del literal 5
z := "Hello, world!" // Se infiere el tipo "string" del literal
z_2 := z             // Inicializar con otra variable toma su tipo de datos
x_2: int = z         // ERROR: z es de tipo "string"
\end{verbatim}
\normalsize

\subsection{Tipos de datos}

El sistema de tipado será \textbf{estático fuerte}: los tipos de datos serán
verificados en tiempo de compilación, y el tipo de datos de una variable será
permanente durante toda la vida de esa variable.

El lenguaje contará con tipos de datos generales, tanto primitivos como
compuestos, y otros específicos al dominio. Los tipos de datos serán:

\begin{description}
	\item[Tipos primitivos:] \texttt{int} (números enteros); \texttt{float}
  	(números decimales, siguiendo IEEE 754); \texttt{bool} (verdadero/falso);
    \texttt{string} (cadenas de texto UTF-8); \texttt{nil} (tipo vacío, su único
    valor válido es \texttt{nil}).
  \item[Tipos de dominio:] \texttt{element}, representa un elemento HTML con su
    nombre de \textit{tag}, atributos, y descendientes. Un archivo fuente que
    retorna una variable de tipo \texttt{element} es considerado un punto de
    entrada, y genera una página en la salida.
  \item[Estructuras:] definidas con la palabra reservada \texttt{struct}, son
    una agrupación de otros valores cada uno con un nombre y tipo de datos.
  \item[Uniones:] definidas con la palabra reservada \texttt{union}, representan
    un tipo que puede tomar uno de varios tipos de datos, junto con información
    del tipo actual de la variable (\textit{tagged union}).
  \item[Opcionales:] definidos a partir de un tipo \texttt{T} como \texttt{T?},
    es un atajo para el patrón común \texttt{union \{ T, nil \}}.
  \item[Enumeraciones:] definidas con la palabra reservada \texttt{enum}, son un
    tipo de datos que toma valores específicos de una lista.
  \item[Listas:] definidas a partir de un tipo \texttt{T} como \texttt{[]T}
    ("lista de \texttt{T}"). Tamaño dinámico.
  \item[Tuplas:] definidas a partir de tipos \texttt{T, U, ...} como
    \texttt{<T, U, ...>}. Tamaño fijo, con elementos de tipos heterogéneos.
  \item[Mapas:] definidos a partir de dos tipos \texttt{K, V} como \texttt{[K]V}
    ("mapa de \texttt{V} indexado por \texttt{K}"). Sintaxis especial para el
    caso \texttt{K = string}, ya que es un uso muy común (al trabajar con JSON).
  \item[Funciones:] definidas a partir de sus tipos de parámetros y de retorno,
    con la palabra reservada \texttt{function}.
\end{description}

Los tipos de datos de usuario se definen con el operador \texttt{is}:

\small
\begin{verbatim}
# Struct usuario con campos int, string y opcional(string)
User is struct {
  id: int,
  username, first_name, last_name: string,
  profile_image_url: string?,
}

# Lista de elementos que pueden tomar uno de cuatro tipos
Article is []union { Paragraph, Image, Quote, CodeBlock }
\end{verbatim}
\normalsize

\subsection{Funciones y componentes}

Las funciones se definen como un tipo de datos de usuario, con la adición de un
cuerpo de función que contiene su código:

\small
\begin{verbatim}
display_name is function(user: User) -> string {
  concat(" ", user.first_name, user.last_name)
}
\end{verbatim}
\normalsize

Toda función debe tener un tipo de retorno, indicado por \texttt{->} siguiendo
su lista de parámetros. Los argumentos de una función se pasan por valor, y no
pueden ser modificados.

Un componente es sencillamente una función que retorna el tipo \texttt{element}:

\small
\begin{verbatim}
user_avatar is function(user: User) -> element {
  div.flex.flex-row.gap-4.items-center {
    // ...
  }
}
\end{verbatim}
\normalsize

La sintaxis de llamada de función es similar a un lenguaje como C, con la
diferencia que si una función no toma parámetros, o todos sus parámetros tienen
un valor por defecto y no se pasa ninguno al llamarla, se puede omitir la lista
de parámetros por completo.

\small
\begin{verbatim}
hello is function(name: string = "world") -> string {
  "Hello, ${name}!"
}

name := display_name(user)
str1 := hello(name)        // "Hello, <users's name>!"
str2 := hello              // "Hello, world!"
\end{verbatim}
\normalsize

Los componentes se definen como cualquier función, pero tienen comportamientos
especiales al llamarlos: se admiten cadenas de texto arbitrarias (sin espacios)
prefijadas con \texttt{.} (cualquier cantidad) o \texttt{\#} (sólo una) entre el
nombre de función y la lista de parámetros. Las primeras se pasan como
\texttt{[]string} al parámetro especial \texttt{class}, y la segunda como
\texttt{string} al parámetro especial \texttt{id}. En caso de no existir estos
parámetros, el uso de su sintaxis especial es un error.

La sintaxis de llamada especial para componentes existe para facilitar la
escritura de código markup:

\small
\begin{verbatim}
// Sin sintaxis especial de llamada
div(class = "flex flex-row gap-4 items-center", id = "myDiv")

// Con sintaxis especial de llamada
div.flex.flex-row.gap-4.items-center#myDiv
\end{verbatim}
\normalsize

\subsection{Composición de funciones}

La construcción que permite utilizar el lenguaje para construir estructuras
HTML complejas de modo sencillo es la \textbf{composición de funciones}. Esto
consiste en pasar una expresión arbitraria como parte de una llamada a función.
La expresión es tratada como un parámetro especial, y es inserta en cada
lugar donde la función tenga la palabra reservada \texttt{compose}:

\small
\begin{verbatim}
twice is function -> <string, string> {
  compose
  compose
}

foo := twice "Hello"    // <"Hello", "Hello">
bar := twice {
  who := "world"
  "Hello ${who}!"
}                       // <"Hello, world!", "Hello, world!">
\end{verbatim}
\normalsize

Un bloque de código pasado para composición captura el \textit{scope} de la
llamada a función; es decir, tiene acceso a las variables del lugar donde se
llama. \textbf{No tiene acceso al \textit{scope} interno de la función.}

\subsection{Expresiones}

El lenguaje será \textbf{basado en expresiones}, es decir, todo en el lenguaje
es una expresión que retorna un valor.

\begin{enumerate}
	\item Una asignación retorna el valor asignado a la variable; esto permite
	  encadenar asignaciones. Las asignaciones tienen comportamiento especial.
	\item Una llamada a función retorna el valor de retorno de la función.
	\item Una expresión condicional retorna el valor de la rama que haya sido
	  tomada en base a la o las condiciones. Todas deben retornar el mismo tipo.
	\item Un \textit{loop} retorna un valor de tipo \texttt{[]T}, donde \texttt{T}
	  es el tipo que retorna la expresión interna.
\end{enumerate}

Un bloque también es una expresión, que consiste de un agrupamiento lineal y
ordenado de otras expresiones. El tipo de retorno de un bloque se decide según
su contexto:

\begin{itemize}
	\item Si el bloque es el cuerpo de una función, su tipo de retorno coincide con
	  el de la función.
	\item En cualquier otro caso, el tipo de retorno se infiere del valor, que
	  sigue las reglas detalladas a continuación.
\end{itemize}

El valor de retorno de un bloque está dado por las siguientes reglas:

\begin{enumerate}
	\item Toda expresión dentro del bloque, \textbf{excepto las asignaciones},
	  agrega su valor de retorno al del bloque.
	\item Si el bloque contiene una única expresión que no es una asignación,
	  retorna directamente su valor.
	\item Si el bloque contiene varias expresiones y todas retornan un mismo tipo
	  \texttt{T}, entonces retorna una lista de tipo \texttt{[]T} con los valores
		de todas sus expresiones (que no sean asignaciones).
	\item Si el bloque contiene varias expresiones y retornan distintos tipos,
	  entonces retorna una tupla con los valores de todas sus expresiones (excepto
		asignaciones).
\end{enumerate}

En el caso del cuerpo de una función, si el valor de retorno no coincide con el
tipo declarado en la función, es un error, excepto en el caso que el bloque
retorne \texttt{[]T} y la función tenga como tipo de retorno una tupla de
\texttt{T} con el tamaño apropiado.

\paragraph{Ejemplos}
\small
\begin{verbatim}
// Retorna string ("Hello, world!")
{
  greeting := "Hello, "
  greeting = greeting + "world!"
  greeting
}

// Retorna []int ([42, 505, 77])
{
  42
  500 + 5
  77
}

// Retorna <int, string> (<42, "Hello, world!">)
{
  (80 + 4) / 2
  who := "world"
  "Hello, ${who}!"
}
\end{verbatim}
\normalsize

\section{Casos de prueba}

Se proponen los siguientes casos de prueba de \textbf{aceptación}:

\begin{enumerate}
	\item Un programa que genere una página simple con texto.
	\item Un programa que genere una página en base a datos de un CSV.
	\item Un programa que genere una página en base a datos de un JSON.
	\item Un programa que defina un componente y lo utilize en una página.
\end{enumerate}

Así como los siguientes casos de prueba de \textbf{rechazo}:

\begin{enumerate}
	\item Un programa malformado.
	\item Un programa que intente operar con tipos de datos incompatibles.
	\item Un programa que asigne una variable que no se utiliza.
	\item Un programa que defina una función que no se utiliza.
	\item Un programa que defina un tipo de datos que no se utiliza.
\end{enumerate}

\section{Ejemplos}

\subsection{Página simple}

Este ejemplo muestra la creación de una página sencilla, utilizando recursos
externos como imágenes y CSS.

\scriptsize
\begin{verbatim}
// "import url" toma un path a recursos en el directorio fuente.
// En compilación, se transforma en url relativa de salida.
my_cool_image := import url "img/cool.png"
tailwind_css  := import url "css/tailwind.css"

html {
  head {
    title "My cool website"
    stylesheet(tailwind_css)
  }
  body.p-8.pb-16 {
    main.max-w-2xl.mx-auto {
      h1.text-2xl.font-bold.mb-4 "Hello, world!"
      p.text-sm.text-black/80.mb-2 {
        "This is some text"
        "and this is some more text"
      }
      img.border.rounded-md(src = my_cool_image)
    }
  }
}
\end{verbatim}
\normalsize

\pagebreak
\subsection{Tabla de CSV}

Este ejemplo muestra la creación de una tabla con datos cargados de un CSV.
Demuestra el uso de \texttt{import csv}, structs, y estructuras de control de
flujo como \textit{loops} y condicionales.

\scriptsize
\begin{verbatim}
Book is struct {
  title: string,
  isbn: string,
  price: float,
  stock: int,
  amount_sold: int,
}

// "import csv" toma un path a recursos en el directorio fuente.
// En compilación, parsea el CSV y retorna los datos con el tipo dado, causando
// un error si los tipos no coinciden.
// En este caso retorna una tupla <[string]string, []Book>
headers, books := import csv of Book "data/csv/books.csv"
tailwind_css   := import url "css/tailwind.css"

html {
  head {
    title "Book inventory"
    stylesheet(tailwind_css)
  }
  body.p-8.pb-16 {
    main.max-w-2xl.mx-auto {
      h1.text-2xl.font-bold.mb-4 "Book Inventory"

      table {
        thead {
          tr.font-bold {
            th.py-1.px-2 headers.title
            th.py-1.px-2 headers.isbn
            th.py-1.px-2 headers.price
            th.py-1.px-2 headers.stock
            th.py-1.px-2 headers.amount_sold
          }
        }
        tbody {
          for book in books {
            tr {
              th.font-semibold.font-italic {
                a.text-underline(href = "/books/${book.isbn}") book.title
              }
              td book.isbn
              td format("$%.2f", book.price)
              td {
                if book.stock == 0 {
                  "Out of stock"
                } else {
                  book.stock "units"
                }
              }
              td { book.amount_sold "units" }
            }
          }
        }
      }
    }
  }
}
\end{verbatim}
\normalsize

\pagebreak
\subsection{Tabla de CSV genérica}

Este ejemplo muestra la definición de un componente para dibujar tablas a partir
de un conjunto de datos arbitrario, y su uso para dibujar una tabla con datos de
un CSV de formato desconocido. Demuestra el uso de componentes, mapas, y de
\texttt{import csv} sin tipo de datos definido.

\scriptsize
\begin{verbatim}
csv_table is function(headers: [string]string, data: [][string]string) -> element {
  fields := for key, header in headers key

  table {
    thead {
      tr.font-bold {
        for key, header in headers th.py-1.px-2 header
      }
    }
    tbody {
      for item in data {
        tr {
          th.font-semibold item[fields[0]]
          for field_name in fields[1:] {
            td item[field_name]
          }
        }
      }
    }
  }
}

// Al no pasarle un tipo de datos, "import csv" retorna una tupla
// <[string]string, [][string]string>, donde las claves de los mapas son
// generadas a partir de los headers.
headers, books := import csv "data/csv/books.csv"
tailwind_css   := import url "css/tailwind.css"

html {
  head {
    title "Book inventory"
    stylesheet(tailwind_css)
  }
  body.p-8.pb-16 {
    main.max-w-2xl.mx-auto {
      h1.text-2xl.font-bold.mb-4 "Book Inventory"

      csv_table(headers, books)
    }
  }
}
\end{verbatim}
\normalsize

\end{document}
